%%This is a very basic article template.
%%There is just one section and two subsections.
\documentclass{article}

\usepackage{url}
\usepackage{caption}
\usepackage{hyperref}
\usepackage[a4paper, margin=2.54cm]{geometry}
\usepackage{wrapfig}
\usepackage{textpos}
\usepackage{tabularx}
\usepackage{color}
\usepackage{ulem}
\usepackage{multirow}
\usepackage{setspace}
\usepackage{graphicx}
\usepackage{color}
\usepackage[]{algorithm2e}
\usepackage{algorithmic}
\usepackage{soul}
\usepackage{soul}
\usepackage{dialogue}

\definecolor{red}{rgb}{1.0, 0.0, 0.0}  %red
\definecolor{TAColor1}{rgb}  {0.0, 0.75,0.0}  %green
\definecolor{TAColor2}{rgb}  {0.0, 0.0, 1.0}  %blue
\definecolor{TAColor3}{rgb}  {1.0, 0.5, 0.0}  %orange
\definecolor{TAColor4}{rgb}  {0.5, 0.5, 0.0}  %dark greenish yellow
\definecolor{TAColor5}{rgb}  {1.0, 0.0, 1.0}  %pink
\definecolor{TAColor6}{rgb}  {0.5, 0.25,0.0}  %brown

\newcommand{\baris}[1]{{\color{red}#1}}
\newcommand{\dogancan}[1]{{\color{TAColor1}#1}}
\newcommand{\onur}[1]{{\color{TAColor2}#1}}
\newcommand{\gursu}[1]{{\color{TAColor3}#1}}
\newcommand{\ainaz}[1]{{\color{TAColor4}#1}}
\newcommand{\ahnaf}[1]{{\color{TAColor6}#1}}

\newif\ifanswers
\answerstrue % comment out to hide answers
\newcommand{\myAnswer}[1]{\ifanswers {\begin{onehalfspacing}
\color{blue}#1
\end{onehalfspacing}} \fi}


%% Definitions
\newcommand{\instructor}{Bar{\i}\c{s} Akg\"un } %need the space at the end
\newcommand{\courseCode}{COMP341 } %need the space at the end
\newcommand{\courseName}{Introduction to Artificial Intelligence}
\newcommand{\assingmentNum}{4}
\newcommand{\responsibleTA}{None}
\newcommand{\dueDate}{June 14 2026, 23:59}
\newcommand{\submissionType}{Blackboard}
\newcommand{\coursePart}{reinforcement learning }

\newcommand{\additionalInfo}{Relevant Programming Session: None}

\newcolumntype{L}{>{\raggedright\arraybackslash}X}%
\newcolumntype{R}{>{\raggedleft\arraybackslash}X}%
\newcolumntype{C}{>{\centering\arraybackslash}X}

\begin{document}

\title{\textbf{Ko\c{c} University \\ \courseCode \\ \courseName \\ Assignment \assingmentNum }}
\date{}

\maketitle

\vspace{-2cm}

\begin{center}
Instructor: \instructor \\
%Responsible TA: \responsibleTA \\
Due Date: \dueDate \\
Submission Through: \submissionType \\
%\additionalInfo \\
\textbf{Make sure you read and understand every part of this document} \\
%\textbf{Important:} Download your submission to make sure it is not corrupted and it has your latest report/code.
%You are only going to be graded by your blackboard submission.
\end{center}

This programming assignment will test your knowledge and your implementation
abilities of what you have learned in the \coursePart part of the class.You are asked to complete
a coding part and answer a few questions about how it runs.  The coding part of the homework will follow the
Berkeley CS188 Spring 2026 Pacman Project P3: Reinforcement Learning at \url{https://inst.eecs.berkeley.edu/~cs188/sp26/projects/proj3/}. 

This homework must be completed individually. Discussion about algorithms, algorithm properties, code structure, and Python is
allowed but group work is not. Coming up with the same approach and talking about the ways of implementation
leads to very similar code which is treated as plagiarism! Furthermore, do not discuss the answers directly as it will lead to
similar sentences which is treated as plagiarism. If you are unsure, you should not discuss. Any academic dishonesty, will not be tolerated.
\textbf{By submitting your assignment, you agree to abide by the Ko\c{c} University codes of conduct.}


\textbf{Important:} Make sure you go over all the submission instructions at the end of this document before submission. Once you upload your submission, 
download it to make sure it is not corrupted and it has your latest documents. You are only going to be graded by your blackboard submission. 

\textbf{Warning:} The solution code for the homeworks can be found online. We are going to compare your submission with these sources. We are also going to 
compare your code to latest and previous submissions of Ko\c{c} students. If your code's similarity level is above a certain threshold, your code will be
scrutinized. If there is strong suspicion of plagiarism, we will take action based on university policies. 

\section*{Programming}
You are going to do the 9 programming questions about \coursePart given in the website.
You are only required to change \textit{valueIterationAgents.py}, \textit{qlearningAgents.py} and \textit{analysis.py}.  If you have any issues with other
parts of the code let your instructor or TA know ASAP, even if you manage to solve your problem.  Use
the data structures in \textit{util.py} for the autograder to work properly. Some of these are really useful as well! If the you think you have the right
answer but the autograder is not giving you any points, try to run it on individual questions (examples on how to do this is given in the website).

The website of the homework and the comments in the code have very useful tips and explanations. Do not skip them!

\subsection*{Hints}
%The questions are fairly straightforward. We have not covered Q4 and Q5 in class, but the website has enough information to implement them. Feel free to
%complete the other questions and get back to these later.

\noindent
\textbf{util.Counter()} \\
The \texttt{util.Counter()} data structure is very useful. It is a modified dictionary, returning 0 for keys that are not in the dictionary. Use this to your
advantage. \\

\noindent
\textbf{Calculating the max and the argmax for Q-Values} \\
I recommend that you write a separate function which returns both the maximum Q-value given a state and the corresponding action (argmax). 
This will be useful in many parts of your code. You can return two values from a function in Python. 
You should do this both for value iteration and q-learning agents. \\

\noindent\textbf{Terminal States} \\
You do not need to calculate anything for the terminal states, do not forget to handle them accordingly! There is a function to test whether a state is
terminal or not. \\

%\noindent
%\textbf{Asynchronous Value Iteration} \\
%You need to iterate over the states in a cyclic manner for the given number of iterations. Use the modulo operator to do this. 



\noindent
This is not a hint but a suggestion; For personal satisfaction, I recommend you add your features from homework 2 and watch pacman learn! 

\section*{Large-Language Model Usage Report}
As mentioned in the syllabus, you are allowed to use Large-Language Models (LLMs) such as ChatGPT, as long as you cite your usage. Towards this end,
you are asked to prepare a short usage report if you used these. This report should include for which questions and how you used LLMs. 
The how part can be achieved by adding a few example prompts and their outputs. I am giving a few examples below:\\

``I used LLMs to get help on Python. I was getting \texttt{unhashable type: \textquotesingle list\textquotesingle} error while 
trying to create an explored set. Here is how I used the LLM: 
\begin{dialogue}
\speak{Me} I am getting an unhashable type: list error while I am trying to populate a list. How can I solve this?
\speak{LLM} The ``unhashable type: list'' error usually occurs when you try to use a list as a key in a dictionary or as an element in a set. 
This is because lists are mutable (changeable) and therefore unhashable, meaning they can't be used as dictionary keys or stored directly in a set. 
To fix this error, check if you're accidentally trying to add a list as an element in a set or as a key in a dictionary. Instead of using a list, 
you can use a tuple (which is immutable) if you need a similar structure but also need it to be hashable.
\end{dialogue}
As a result, I have changed my state representation from a list to a tuple and fixed my issue.'' \\

You can also describe it as a narrative without all the interaction: \\
``I used LLMs to get an implementation of the depth first algorithm. 
My prompt was ``...." followed by the incomplete code I needed to complete. The response was not compatible with the code base so 
I additionally sent this ``...". Finally, I modified line 7 so that the dictionary has the correct values."

It does not have to be code, you can use it for math as well:\\
``I forgot how the quadratic formula and this is how I used ChatGPT
\begin{dialogue}
\speak{Me} Can you explain how to solve a quadratic equation using the quadratic formula, with an example?
\speak{LLM} To solve a quadratic equation of the form $ax^2 + bx + c = 0$ using the quadratic formula, you apply the formula $x = (-b ± \sqrt{b²-4ac}) / (2a)$. 
For example, for the equation $x^2 - 3x - 4 = 0$, we have $x_1 = $ \ldots
\end{dialogue}

Prepare and send a report with the aforementioned information. If the report looks like it is going to be too long, tell us where you used it and only
provide a few examples. The format of the report is not too critical but it needs to clearly show how you used LLMs.


\section*{Submission}

You are going to submit a compressed archive through the blackboard site. 
This file should only contain \textit{valueIterationAgents.py}, \textit{qlearningAgents.py}, \textit{analysis.py} and optionally \textit{report.pdf} if you used LLMs. 
Other files will be deleted and/or overwritten. 


\subsection*{Submission Instructions} 
\begin{itemize}
 \item You are going to submit a \textbf{single} compressed archive through the blackboard site. 
 The file can have \textit{zip}, \textit{rar},  \textit{tar}, \textit{tar.gz} or \textit{7z} format. 
 If you submit files individually they may get lost.
 \item You are fine as long as the compressed archive has the required files within 4 folder levels.
 \item Code that does not run (e.g. due to syntax errors), that does not terminate (e.g. due to infinite loops) or that blows up memory will not get any
 points.
 \item \textbf{Important:} Download your submission to make sure it is not corrupted and it has your latest code. 
You are only going to be graded by your blackboard submission.
\end{itemize}

\end{document}
